\chapter{Perguntas e respostas ao perito}
\label{anx:sessoes}

{\footnotesize
\begin{longtable}{@{}L{2.6cm}VL{5.3cm}VL{7.0cm}@{}}
  \caption{Perguntas e respostas ao perito}
  \label{tab:perguntas-perito}\\
  \toprule
  \textbf{Tema} & \textbf{Pergunta (síntese)}
    & \textbf{Resposta do perito (síntese)} \\
  \midrule
  \endfirsthead
  \toprule
  \textbf{Tema} & \textbf{Pergunta (síntese)}
    & \textbf{Resposta do perito (síntese)} \\
  \midrule
  \endhead
  \bottomrule
  \endfoot
  Falha do veículo
    & Que estado deve assumir o veículo em caso de falha?
    & \enquote{STOP imediato, sem dúvida.} \\ \linhasep
  Relevância dos sensores
    & Que informação dos sensores é mais importante (velocidade, distância
      ao obstáculo, TTC, semáforos, meteorologia, sinais de trânsito)?
    & Os dados indicados já são suficientes. Todos são importantes: cada
      sensor é necessário na situação específica para a qual contribui. \\ \linhasep
  Criticidade e TTC
    & Como distinguir uma situação normal de uma crítica, usando o TTC para
      decidir entre desviar e travar?
    & A criticidade deve resultar de uma equação a desenvolver pela equipa,
      com base na velocidade atual, na distância ao objeto e na capacidade de
      travagem. Existe um \enquote{ponto de não-retorno} a partir do qual a colisão é
      inevitável. \enquote{Nunca considerem arranque e travagem instantâneos.} \\ \linhasep
  Número de leituras
    & Quantas leituras considerar antes de decidir?
    & Depende do resultado da equação. Numa situação crítica mas ainda
      passível de manobra, a decisão é tomada logo à primeira leitura (por
      exemplo, mudança de via). Numa situação normal, pode fazer-se uma
      segunda leitura para excluir um valor atípico (falso positivo). \\ \linhasep
  Prioridades e desempate
    & Que níveis de prioridade definir e como desempatar situações com a
      mesma prioridade?
    & Usar quatro níveis: baixo, médio, alto e crítico. Para desempatar,
      escolher aleatoriamente entre as possibilidades ou, se estiverem
      ordenadas numa lista, escolher a primeira. \\ \linhasep
  Memória de trabalho
    & Durante quanto tempo guardar factos como a presença de um peão ou de
      um obstáculo?
    & Quando o evento sai do campo de deteção, passa a ser um evento passado
      e pode ser esquecido. Não é necessário memorizar eventos: uma alteração
      por sinal de trânsito é uma mudança de estado, não uma posição a
      memorizar. \\ \linhasep
  Dados contraditórios
    & Como deve o veículo comportar-se quando os sensores dão informação
      contraditória?
    & Sem uma terceira fonte, aplica-se o critério de desempate. Se o nível
      for alto ou crítico, \enquote{é melhor o STOP}. \\ \linhasep
  Cenários
    & Que cenários apresentar?
    & Todos os discutidos na reunião (S1), podendo acrescentar-se a
      passadeira com peão. Sugere ainda permitir que o utilizador clique na
      estrada para fazer surgir um elemento fora de contexto (por exemplo,
      carga caída de outra via). \\ \linhasep
  Explicações
    & Como apresentar as explicações \emph{Why} e \emph{Why Not}: pelas
      regras ou por uma explicação?
    & As regras são obrigatórias e a informação apresentada deve estar
      alinhada com elas. \\ \linhasep
  Melhorias
    & Que melhorias sugere?
    & Considera difícil sugerir melhorias a um sistema que ainda não existe,
      mas avalia como correto o levantamento de requisitos e o plano de
      trabalhos. Devemos voltar a fazer esta questão numa fase mais avançada
      do projeto. \\
\end{longtable}
}
